\documentclass[10pt,oneside]{article}

\usepackage[letterpaper,top=0.82in,bottom=0.82in,left=1.02in,right=1.02in]{geometry}
\usepackage[T1]{fontenc}
\usepackage[sc]{mathpazo}
\usepackage{microtype}
\usepackage{amsmath,amssymb,amsthm,mathtools}
\usepackage{bm}
\usepackage{enumitem}
\usepackage{titlesec}
\usepackage{fancyhdr}
\usepackage[hidelinks]{hyperref}
\usepackage{bookmark}
\usepackage{xcolor}

\hypersetup{
  pdftitle={Motion and Difference},
  pdfauthor={Anthony Vito Coppa},
  pdfsubject={Projective arithmetic and Lorentzian gauge compatibility},
  pdfkeywords={projective dynamics, Fibonacci matrix, hyperbolic geometry,
    Lorentzian gauge field, holonomy, Maxwell equations}
}

\definecolor{ink}{RGB}{18,18,18}
\color{ink}

\setlength{\parindent}{1.15em}
\setlength{\parskip}{0pt}
\setlength{\columnsep}{0.25in}
\linespread{1.035}
\allowdisplaybreaks[2]

\titleformat{\section}
  {\normalfont\small\scshape\centering}
  {\Roman{section}.}{0.55em}{}
\titleformat{\subsection}
  {\normalfont\normalsize\itshape}
  {\thesubsection.}{0.45em}{}
\titlespacing*{\section}{0pt}{2.2ex plus .7ex minus .2ex}{1.2ex}
\titlespacing*{\subsection}{0pt}{1.7ex plus .5ex minus .2ex}{0.7ex}

\newtheoremstyle{plaincompact}
  {1.1ex}{1.1ex}{\itshape}{}{}{.}{0.55em}{}
\theoremstyle{plaincompact}
\newtheorem{theorem}{Theorem}
\newtheorem{proposition}[theorem]{Proposition}
\newtheorem{lemma}[theorem]{Lemma}
\newtheorem{corollary}[theorem]{Corollary}

\newtheoremstyle{definitioncompact}
  {1.1ex}{1.1ex}{\normalfont}{}{}{.}{0.55em}{}
\theoremstyle{definitioncompact}
\newtheorem{definition}[theorem]{Definition}
\newtheorem{remark}[theorem]{Remark}

\renewcommand{\qedsymbol}{$\square$}
\newcommand{\R}{\mathbb R}
\newcommand{\C}{\mathbb C}
\newcommand{\Hh}{\mathbb H}
\newcommand{\RP}{\mathbb {RP}}
\newcommand{\Fib}{F}
\newcommand{\id}{\operatorname{id}}
\newcommand{\Fix}{\operatorname{Fix}}
\newcommand{\Eq}{\operatorname{Eq}}
\newcommand{\Hol}{\operatorname{Hol}}
\newcommand{\vol}{\operatorname{vol}}
\newcommand{\Tr}{\operatorname{tr}}
\newcommand{\Ad}{\mathcal A}
\newcommand{\cF}{\mathcal F}
\newcommand{\cT}{\mathcal T}
\newcommand{\cU}{\mathcal U}
\newcommand{\dd}{\,\mathrm d}

\pagestyle{fancy}
\fancyhf{}
\fancyhead[L]{\footnotesize\scshape Motion and Difference}
\fancyhead[R]{\footnotesize A. V. Coppa}
\fancyfoot[C]{\footnotesize\thepage}
\renewcommand{\headrulewidth}{0.35pt}
\setlength{\headheight}{13pt}
\fancypagestyle{first}{
  \fancyhf{}
  \fancyfoot[C]{\footnotesize\thepage}
  \renewcommand{\headrulewidth}{0pt}
}

\begin{document}

\thispagestyle{first}
\begin{center}
  \vspace*{0.42in}
  {\fontsize{23}{27}\selectfont\scshape Motion and Difference\par}
  \vspace{0.34in}
  {\large Anthony Vito Coppa\par}
  \vspace{0.08in}
  {\small Independent Researcher\par}
  \vspace{0.08in}
  {\small July 17, 2026\par}
  \vspace{0.08in}
  {\small\href{https://doi.org/10.5281/zenodo.22780641}{doi:10.5281/zenodo.22780641}\par}
\end{center}

\vspace{0.34in}

\begin{abstract}
\noindent
We construct one explicit mathematical realization in which a two-dimensional
carrier supports complex phase, square symmetry, projective arithmetic
reflection, Fibonacci refinement, and a Lorentzian $SO(2)$ gauge field.  The
elementary factorization
$Q=M_JC_1=\left(\begin{smallmatrix}1&1\\1&0\end{smallmatrix}\right)$
turns alternating projective involutions into the exact matrix record of the
Fibonacci sequence.  Its two fixed points form a maximal equalizer in
$PGL(2,\R)$, determine a hyperbolic glide axis, and recover the sharp
Diophantine scale $\sqrt5$.  On a Lorentzian four-manifold, the same phase
generator $C_1$ carries a global $SO(2)$ connection.  A nonnegative gluing
functional detects exact compatibility of local bundle data.  An explicit
one-parameter family of connections keeps a nonzero source-free Maxwell
curvature fixed while varying relative holonomy independently; a pair of
endpoint-sharing timelike paths also has unequal proper time.  The completed
zeta reflection is placed on the same arithmetic involution, with no asserted
spectral, causal, or functional identification among the resulting scalar and
geometric quantities.  The main result is therefore a compatibility theorem:
the declared structures coexist exactly in a single realization, while the
boundaries between them remain explicit.
\end{abstract}

\vspace{0.12in}
\noindent\textit{Keywords:} projective dynamics; Fibonacci matrix; hyperbolic
geometry; Diophantine approximation; Lorentzian gauge field; holonomy;
Maxwell equations.

\vspace{0.18in}
\hrule
\vspace{0.10in}
\noindent\footnotesize
\textit{Classification:} 11B39, 11J70, 20H10, 53C29, 58J28, 83C50.
\normalsize

\section{Statement and scope}

The object of this paper is narrow.  We ask whether several exact structures,
normally written in separate mathematical languages, can be carried without
contradiction by one explicit model.  The answer is affirmative.  The model
contains a common two-dimensional carrier, but it does not collapse the roles
placed on that carrier.

The arithmetic core begins with two involutions on the real projective line.
One is induced by a complex quarter-turn; the other is induced by the affine
reflection $x\mapsto 1-x$.  Their alternating composition is represented by
the classical Fibonacci matrix.  The same reflection is also the argument
symmetry of the completed Riemann zeta function.  These statements meet at a
precisely declared two-point set and nowhere else.

The geometric core begins with an oriented rank-two Euclidean bundle over a
Lorentzian four-manifold.  Its $SO(2)$ connection uses the same infinitesimal
quarter-turn generator.  Exact gluing is detected by a positive functional.
In an explicit product spacetime, the connection has nonzero source-free
Maxwell curvature, independently variable holonomy, and timelike paths of
unequal proper time.

The classical ingredients concerning Fibonacci numbers, continued fractions,
Hurwitz approximation, hyperbolic isometries, the completed zeta function,
connections, and Maxwell fields are standard.  The point proved here is their
exact placement in a single bounded construction, centered on the factorization
$Q=M_JC_1$ and completed by an explicit Lorentzian witness.

\begin{theorem}[Simultaneous realization]\label{thm:main}
There exists a tuple
\[
\mathfrak A=(V,C_1,\Pi,M_J,Q,J,\xi,M,g,E,\Ad,q,\gamma_0,\gamma_1)
\]
with the following properties.
\begin{enumerate}[label=(\roman*),leftmargin=2.25em,itemsep=0.25ex,topsep=0.5ex]
\item $\langle C_1,\Pi\rangle\cong D_4$ on $V=\R^2$ and
      $\mathbb P(D_4)\cong V_4$;
\item $Q=M_JC_1$, $Q^2=Q+I$, and the alternating projective iterate has
      Fibonacci coefficients;
\item the equalizer of the two projective involutions is maximal and consists
      exactly of the two roots of $s^2-s-1$;
\item $Q$ acts on the upper half-plane as a glide reflection of displacement
      $2\log\Phi$, and the Fibonacci convergents attain the sharp Hurwitz scale
      asymptotically;
\item the completed zeta function obeys $\xi\circ J=\xi$, including equality
      on the two refinement roots;
\item $E\to M$ carries an $SO(2)$ connection with nonzero source-free Maxwell
      curvature and exact local gluing;
\item within a one-parameter family, curvature is fixed while relative holonomy
      is arbitrarily prescribable in $SO(2)$;
\item two timelike paths with common endpoints have unequal proper time.
\end{enumerate}
No equality or functional dependence is asserted among the zeta remainder,
Maxwell curvature, relative holonomy, proper-time difference, and gluing
functional.
\end{theorem}

The proof occupies Sections~II--VIII.  Each interface is typed before it is
used, and the nonidentifications are collected in Section~IX.

\section{The carrier and the alternating factorization}

Let $V=\R^2$ with its standard Euclidean inner product.  We use the same real
space as $\C$ for phase and as $\mathbb P(V)\cong\RP^1$ for projective action.
Define
\begin{equation}\label{eq:generators}
C_1=\begin{pmatrix}0&-1\\1&0\end{pmatrix},\qquad
\Pi=\begin{pmatrix}1&0\\0&-1\end{pmatrix},\qquad
M_J=\begin{pmatrix}-1&1\\0&1\end{pmatrix}.
\end{equation}
Then
\begin{equation}\label{eq:d4relations}
C_1^2=-I,\qquad C_1^4=I,\qquad
\Pi^2=I,\qquad \Pi C_1\Pi=C_1^{-1},\qquad M_J^2=I.
\end{equation}
Thus $\langle C_1,\Pi\rangle\cong D_4$, the order-eight symmetry group of the
square.  Under $V\cong\C$, multiplication by $i$ is $C_1$, and
\begin{equation}\label{eq:exp}
e^{\theta C_1}=I\cos\theta+C_1\sin\theta,
\qquad e^{\pi C_1}=-I,\qquad e^{2\pi C_1}=I.
\end{equation}
The projective image has kernel $\{I,-I\}$, hence
\begin{equation}\label{eq:projectived4}
\mathbb P(D_4)\cong D_4/\{\pm I\}\cong V_4.
\end{equation}

Let $\mathbb P(A)$ denote the projective map induced by an invertible real
matrix $A$.  Put
\begin{equation}\label{eq:projmaps}
\sigma=\mathbb P(C_1),\qquad J_{\mathbb P}=\mathbb P(M_J).
\end{equation}
On the affine chart $\iota(x)=[x:1]$, these are
\begin{equation}\label{eq:affineinvolutions}
\sigma_{\mathrm{aff}}(x)=-\frac1x,\qquad
J_{\mathrm{aff}}(x)=1-x.
\end{equation}
Both projective maps are involutions.  They are distinct because the signs of
$\det C_1=1$ and $\det M_J=-1$ cannot be changed by rescaling a real matrix
representative.

The central factorization is
\begin{equation}\label{eq:factorization}
\boxed{
Q=M_JC_1=
\begin{pmatrix}1&1\\1&0\end{pmatrix}.}
\end{equation}
Its characteristic polynomial is $\lambda^2-\lambda-1$.  Cayley-Hamilton gives
\begin{equation}\label{eq:refinementlaw}
Q^2=Q+I,\qquad Q^{-1}=Q-I=C_1^{-1}M_J.
\end{equation}
Defining $\cT=\mathbb P(Q)$, we obtain
\begin{equation}\label{eq:orderedcomposition}
\cT=J_{\mathbb P}\circ\sigma,\qquad
\cT^{-1}=\sigma\circ J_{\mathbb P}.
\end{equation}
On the affine chart,
\begin{equation}\label{eq:T}
T(x)=1+\frac1x,\qquad T^{-1}(x)=\frac1{x-1}.
\end{equation}
The order matters.  Directly,
\begin{equation}\label{eq:commutator}
[Q,C_1]=QC_1-C_1Q=
\begin{pmatrix}2&-1\\-1&-2\end{pmatrix}\ne0.
\end{equation}

\subsection{The maximal meeting}

Let $P(s)=s^2-s-1$ and define the affine reflection
\begin{equation}\label{eq:J}
J(s)=1-s.
\end{equation}
The polynomial is invariant under $J$:
\begin{equation}\label{eq:Pinvariance}
P\circ J=P.
\end{equation}
Its roots are
\begin{equation}\label{eq:roots}
\Phi=\frac{1+\sqrt5}{2},\qquad
J(\Phi)=\frac{1-\sqrt5}{2}=-\frac1\Phi.
\end{equation}
Equivalently, a whole-remainder proportion $W=L+R$ and
$W/L=L/R=x$ yields $x=T(x)$, hence $P(x)=0$.  In the centered coordinate
$w=s-1/2$, the same equation is $w^2=5/4$.

\begin{proposition}[Maximal equalizer]\label{prop:equalizer}
Let $\mathcal R_\Phi=\{\Phi,J(\Phi)\}$.  Then
\begin{equation}\label{eq:equalizer}
\Eq(J_{\mathbb P},\sigma)=\Fix(\cT)=\Fix(\cT^{-1})
=\{[x:1]:x\in\mathcal R_\Phi\}.
\end{equation}
This two-point set is the largest possible equalizer of two distinct elements
of $PGL(2,\R)$.
\end{proposition}

\begin{proof}
Since $J_{\mathbb P}$ is an involution,
$J_{\mathbb P}(p)=\sigma(p)$ if and only if
$(J_{\mathbb P}\circ\sigma)(p)=p$.  Neither $0$ nor $\infty$ is fixed by
$\cT$.  On the affine chart, $T(x)=x$ if and only if $P(x)=0$, giving the two
points in Eq.~\eqref{eq:equalizer}.  The inverse has the same fixed points.
For distinct $f,g\in PGL(2,\R)$,
$\Eq(f,g)=\Fix(g^{-1}f)$.  A nonidentity real projective transformation has at
most two real eigenlines, hence at most two fixed points in $\RP^1$.
\end{proof}

\begin{corollary}\label{cor:dinfty}
The group generated by the two projective involutions is infinite dihedral:
\begin{equation}\label{eq:dinfty}
\langle J_{\mathbb P},\sigma\rangle\cong D_\infty.
\end{equation}
\end{corollary}

\begin{proof}
The eigenvalues of $Q$ are $\Phi$ and $-\Phi^{-1}$.  Their ratio has modulus
$\Phi^2\ne1$, so $\cT=\mathbb P(Q)$ has infinite order.  Two involutions whose
product has infinite order generate $D_\infty$.
\end{proof}

\section{Fibonacci dynamics and the hyperbolic axis}

Let $\Fib_0=0$, $\Fib_1=1$, and
$\Fib_{n+1}=\Fib_n+\Fib_{n-1}$ for $n\ge1$.

\begin{proposition}[Fibonacci powers]\label{prop:fibpowers}
For every $n\ge1$,
\begin{equation}\label{eq:fibpowers}
Q^n=\Fib_nQ+\Fib_{n-1}I
=\begin{pmatrix}
\Fib_{n+1}&\Fib_n\\
\Fib_n&\Fib_{n-1}
\end{pmatrix}.
\end{equation}
Consequently,
\begin{equation}\label{eq:projectivefib}
(J_{\mathbb P}\circ\sigma)^n=\mathbb P(Q^n),
\qquad
\cT^n(\infty)=\frac{\Fib_{n+1}}{\Fib_n}.
\end{equation}
\end{proposition}

\begin{proof}
The assertion holds for $n=1$.  If
$Q^n=\Fib_nQ+\Fib_{n-1}I$, then Eq.~\eqref{eq:refinementlaw} gives
\[
Q^{n+1}=\Fib_n(Q+I)+\Fib_{n-1}Q
=\Fib_{n+1}Q+\Fib_nI.
\]
The matrix formula and projective statements follow.
\end{proof}

The ratios in Eq.~\eqref{eq:projectivefib} are the finite continued fractions
$[1;1,\ldots,1]$ with $n$ entries.  The exact Binet normalization is
\begin{equation}\label{eq:binet}
\Fib_n=\frac{\Phi^n-J(\Phi)^n}{\Phi-J(\Phi)}
=\frac{\Phi^n-J(\Phi)^n}{\sqrt5},
\end{equation}
and
\begin{equation}\label{eq:exacterror}
\Fib_{n+1}-\Phi\Fib_n=J(\Phi)^n.
\end{equation}
It follows that
\begin{equation}\label{eq:scalederror}
\sqrt5\,\Fib_n^2
\left|\Phi-\frac{\Fib_{n+1}}{\Fib_n}\right|
=1-(-1)^n\Phi^{-2n},
\end{equation}
so
\begin{equation}\label{eq:hurwitzlimit}
\lim_{n\to\infty}\sqrt5\,\Fib_n^2
\left|\Phi-\frac{\Fib_{n+1}}{\Fib_n}\right|=1.
\end{equation}
Hurwitz's theorem~\cite{HardyWright,Khinchin} states that every irrational
$\alpha$ has infinitely many
rational approximants $p/q$ with
\begin{equation}\label{eq:hurwitz}
\left|\alpha-\frac pq\right|<\frac1{\sqrt5\,q^2},
\end{equation}
and that the constant is optimal uniformly over irrational $\alpha$.
Equation~\eqref{eq:hurwitzlimit} places the Fibonacci convergents to $\Phi$ at
that sharp asymptotic scale.  The same $\sqrt5$ is also the exact separation
\begin{equation}\label{eq:rootgap}
\Phi-J(\Phi)=\sqrt5
\end{equation}
of the two refinement branches.  These are two declared mathematical offices
of one constant, not a causal relation.

Taking determinants in Eq.~\eqref{eq:fibpowers} yields Cassini's identity,
\begin{equation}\label{eq:cassini}
\Fib_{n+1}\Fib_{n-1}-\Fib_n^2=(-1)^n.
\end{equation}
Thus even projective iterates preserve orientation and odd iterates reverse it.

\subsection{Upper-half-plane realization}

For $A=\left(\begin{smallmatrix}a&b\\c&d\end{smallmatrix}\right)\in
GL(2,\R)$, define the standard extended action on $\Hh$ by
Refs.~\cite{Beardon,Lang}:
\begin{equation}\label{eq:extendedaction}
\rho_A(z)=
\begin{cases}
(az+b)/(cz+d),&\det A>0,\\[3pt]
(a\overline z+b)/(c\overline z+d),&\det A<0.
\end{cases}
\end{equation}
Then
\begin{equation}\label{eq:hyperbolicmaps}
\rho_{C_1}(z)=-\frac1z,\qquad
\rho_{M_J}(z)=1-\overline z,\qquad
\rho_Q(z)=1+\frac1{\overline z}.
\end{equation}
The first is rotation by $\pi$ about $i$.  The second is reflection in
$\operatorname{Re}z=1/2$.  The third preserves the geodesic with endpoints
$J(\Phi)$ and $\Phi$,
\begin{equation}\label{eq:axis}
\left|z-\frac12\right|=\frac{\sqrt5}{2},\qquad \operatorname{Im}z>0.
\end{equation}

\begin{proposition}[Glide displacement]\label{prop:glide}
The map $\rho_Q$ is an orientation-reversing glide reflection along the axis
\eqref{eq:axis}, with glide displacement $2\log\Phi$.  Its square is a
hyperbolic translation of length $4\log\Phi$.
\end{proposition}

\begin{proof}
The boundary fixed points solve $x=1+1/x$, hence are the two points of
$\mathcal R_\Phi$.  The map has no fixed point in $\Hh$, so it is a glide
reflection along their geodesic.  Moreover,
\[
Q^2=\begin{pmatrix}2&1\\1&1\end{pmatrix}\in SL(2,\R),
\qquad \Tr(Q^2)=3,
\]
with eigenvalues $\Phi^2$ and $\Phi^{-2}$.  Its translation length is
\begin{equation}\label{eq:translationlength}
2\operatorname{arccosh}(3/2)=4\log\Phi.
\end{equation}
A glide reflection squares to translation by twice its displacement.
\end{proof}

Equations~\eqref{eq:fibpowers} and \eqref{eq:hyperbolicmaps} therefore describe
one discrete orbit in three compatible readings: integer matrix coefficients,
continued-fraction refinement, and alternating translation-reflection along a
fixed hyperbolic axis.

\section{The completed-zeta reflection}

Let $\zeta_{\mathrm R}$ denote the analytically continued Riemann zeta function
and define
\begin{equation}\label{eq:xi}
\xi(s)=\frac12s(s-1)\pi^{-s/2}
\Gamma\!\left(\frac s2\right)\zeta_{\mathrm R}(s).
\end{equation}
The function $\xi$ is entire and satisfies the classical functional
equation~\cite{Titchmarsh}
\begin{equation}\label{eq:xifunctional}
\xi(s)=\xi(1-s),
\end{equation}
or, using Eq.~\eqref{eq:J},
\begin{equation}\label{eq:xiJ}
\xi\circ J=\xi.
\end{equation}
Since the refinement roots form a $J$-pair,
\begin{equation}\label{eq:xiroots}
\xi(\Phi)=\xi(J(\Phi)).
\end{equation}
Also, because $\sigma_{\mathrm{aff}}$ exchanges these same two points,
\begin{equation}\label{eq:restrictedxi}
\left.(\xi\circ\sigma_{\mathrm{aff}})\right|_{\mathcal R_\Phi}
=\left.\xi\right|_{\mathcal R_\Phi}.
\end{equation}
Equation~\eqref{eq:restrictedxi} is a two-point restriction, not a global
identity on $\R^\times$.

For the antiholomorphic reflection $J_{\Hh}(s)=1-\overline s$, Schwarz
reflection and Eq.~\eqref{eq:xifunctional} give
\begin{equation}\label{eq:modxi}
\xi(J_{\Hh}(s))=\overline{\xi(s)},
\qquad |\xi(J_{\Hh}(s))|=|\xi(s)|.
\end{equation}
This is the standard symmetry about the critical line.  Nothing here implies a
claim about the location of zeta zeros.

For later separation of types, define only the scalar remainder
\begin{equation}\label{eq:epsilon}
\epsilon_\Phi=\zeta_{\mathrm R}(\Phi)-\sqrt5.
\end{equation}
Then $\zeta_{\mathrm R}(\Phi)=\sqrt5+\epsilon_\Phi$ is an identity by
definition.  No magnitude, geometric meaning, or coupling is assigned to
$\epsilon_\Phi$.

\section{Lorentzian bundle structure and exact gluing}

Let $(M,g)$ be a time-oriented Lorentzian four-manifold and let $h$ be an
auxiliary Riemannian metric on $M$.  Let $\pi:E\to M$ be an oriented rank-two
real vector bundle with structure group $SO(2)$ and positive-definite fiber
metric $h_E$.  On a finite trivializing cover $\{U_\alpha\}$, take local
section representatives $q_\alpha:U_\alpha\to V$, transition maps
$\cU_{\alpha\beta}:U_{\alpha\beta}\to SO(2)$, and local connection one-forms
$\Ad_\alpha\in\Omega^1(U_\alpha)$; see
Refs.~\cite{KobayashiNomizu,Nakahara,BaezMuniain}.

Normalize
\begin{equation}\label{eq:normalization}
\cU_{\alpha\alpha}=I,\qquad
\cU_{\beta\alpha}=\cU_{\alpha\beta}^{-1},
\end{equation}
and choose overlaps on which
\begin{equation}\label{eq:lifts}
\cU_{\alpha\beta}=e^{\chi_{\alpha\beta}C_1}
\end{equation}
for smooth real lifts $\chi_{\alpha\beta}$.  Different lifts differ locally by
$2\pi\mathbb Z$, so $d\chi_{\alpha\beta}$ is lift-independent.  Define
\begin{equation}\label{eq:covariantderivative}
D_\alpha q_\alpha=dq_\alpha+\Ad_\alpha C_1q_\alpha.
\end{equation}
The connection law
\begin{equation}\label{eq:connectionlaw}
\Ad_\beta=\Ad_\alpha-d\chi_{\alpha\beta}
\end{equation}
makes $F_\alpha=d\Ad_\alpha$ global.  Because
$\mathfrak{so}(2)=\R C_1$, the matrix-valued connection and curvature are
\begin{equation}\label{eq:matrixcurvature}
\widehat\Ad_\alpha=\Ad_\alpha C_1,\qquad
\Omega_\alpha=d\widehat\Ad_\alpha+
\widehat\Ad_\alpha\wedge\widehat\Ad_\alpha=F_\alpha C_1.
\end{equation}
The Abelian quadratic term vanishes.  Thus $C_1$ is at once the complex phase
generator, the quarter-turn in $D_4$, the infinitesimal generator of parallel
transport, and the Lie-algebra value of curvature.  This is shared carrier
structure.  It is not equality of the scalar quantities multiplying $C_1$.

The arithmetic reflection remains outside the gauge group.  Indeed,
\begin{equation}\label{eq:MJnotorthogonal}
M_J^{T}M_J=\begin{pmatrix}1&-1\\-1&2\end{pmatrix}\ne I,
\qquad M_J\notin O(2).
\end{equation}
Hence $M_J$ is neither an $SO(2)$ structure transformation nor a fiber
isometry.

\subsection{A zero-detecting functional}

Define the local defects
\begin{align}
\Delta^q_{\alpha\beta}
  &=q_\beta-\cU_{\alpha\beta}q_\alpha,
  \label{eq:qdefect}\\
\Delta^{\Ad}_{\alpha\beta}
  &=\Ad_\beta-\Ad_\alpha+d\chi_{\alpha\beta},
  \label{eq:Adefect}\\
\Delta^{\cU}_{\alpha\beta\gamma}
  &=\cU_{\alpha\beta}\cU_{\beta\gamma}
    \cU_{\gamma\alpha}-I.
  \label{eq:Udefect}
\end{align}
Assume continuity and square-integrability with respect to $h$ and $h_E$.
For fixed $\kappa_{\Ad},\kappa_{\cU}>0$, set
\begin{align}
\cF_q
  &=\sum_{\alpha<\beta}\int_{U_{\alpha\beta}}
    \|\Delta^q_{\alpha\beta}\|_{h_E}^2\dd\vol_h,
  \label{eq:Fq}\\
\cF_{\Ad}
  &=\sum_{\alpha<\beta}\int_{U_{\alpha\beta}}
    \|\Delta^{\Ad}_{\alpha\beta}\|_{h}^2\dd\vol_h,
  \label{eq:FA}\\
\cF_{\cU}
  &=\sum_{\substack{\alpha<\beta<\gamma\\
      U_{\alpha\beta\gamma}\ne\varnothing}}
    \int_{U_{\alpha\beta\gamma}}
    \|\Delta^{\cU}_{\alpha\beta\gamma}\|_{\mathrm{HS}}^2\dd\vol_h,
  \label{eq:FU}\\
\cF_\Lambda
  &=\cF_q+\kappa_{\Ad}\cF_{\Ad}
    +\kappa_{\cU}\cF_{\cU}.
  \label{eq:Flambda}
\end{align}

\begin{proposition}[Exact gluing]\label{prop:gluing}
Under the stated assumptions,
\begin{equation}\label{eq:gluingiff}
\cF_\Lambda=0
\quad\Longleftrightarrow\quad
\Delta^q_{\alpha\beta}=0,
\quad\Delta^{\Ad}_{\alpha\beta}=0,
\quad\Delta^{\cU}_{\alpha\beta\gamma}=0
\end{equation}
on every relevant overlap.
\end{proposition}

\begin{proof}
Every integrand in Eq.~\eqref{eq:Flambda} is nonnegative.  If all defects
vanish, the functional vanishes.  Conversely, a zero finite sum with positive
coefficients forces each integral to vanish.  If a continuous squared norm
were positive at one point, it would remain positive on a nonempty open
neighborhood of positive $h$-volume, a contradiction.
\end{proof}

The auxiliary metric changes the numerical value of $\cF_\Lambda$ but not its
zero locus.

\subsection{Mirror covariance}

Under
\begin{equation}\label{eq:mirroraction}
q_\alpha\mapsto\Pi q_\alpha,\qquad
\Ad_\alpha\mapsto-\Ad_\alpha,\qquad
\chi_{\alpha\beta}\mapsto-\chi_{\alpha\beta},\qquad
\cU_{\alpha\beta}\mapsto\Pi\cU_{\alpha\beta}\Pi,
\end{equation}
the defects transform as
\begin{equation}\label{eq:mirrordefects}
\Delta^q\mapsto\Pi\Delta^q,\qquad
\Delta^{\Ad}\mapsto-\Delta^{\Ad},
\qquad \Delta^{\cU}\mapsto\Pi\Delta^{\cU}\Pi.
\end{equation}
Orthogonality of $\Pi$ therefore gives
\begin{equation}\label{eq:mirrorfunctional}
\cF_\Lambda(\Pi q,-\Ad,\Pi\cU\Pi;h,h_E)
=\cF_\Lambda(q,\Ad,\cU;h,h_E).
\end{equation}
Curvature changes sign, $F\mapsto-F$, preserving the source-free Maxwell
equations $dF=0$ and $d\star_gF=0$.

\section{Relative phase, proper time, and independent tuning}

For a piecewise-smooth path $\gamma:[0,1]\to M$, parallel transport solves
\begin{equation}\label{eq:paralleltransport}
\dot v+\Ad(\dot\gamma)C_1v=0.
\end{equation}
In a global Abelian trivialization,
\begin{equation}\label{eq:pathtransport}
P_\Ad(\gamma)=\exp\!\left(-C_1\int_\gamma\Ad\right).
\end{equation}
For paths $\gamma_1,\gamma_2:p\to q$, let
$\ell_{12}=\gamma_1*\overline{\gamma_2}$.  Then
\begin{equation}\label{eq:relativeholonomy}
\Hol_\Ad(\ell_{12})=P_\Ad(\gamma_2)^{-1}P_\Ad(\gamma_1),
\qquad
\Delta_{\mathrm{phase}}(\gamma_1,\gamma_2)
=\Hol_\Ad(\ell_{12})-I.
\end{equation}
If the paths are timelike, define
\begin{equation}\label{eq:propertime}
\tau[\gamma]=\frac1c\int
\sqrt{-g(\dot\gamma,\dot\gamma)}\dd\lambda,
\qquad
\Delta_{\mathrm{time}}(\gamma_1,\gamma_2)
=\tau[\gamma_1]-\tau[\gamma_2].
\end{equation}
The two differences have different codomains:
$\Delta_{\mathrm{phase}}\in\operatorname{End}(E_p)$ and
$\Delta_{\mathrm{time}}\in\R$.

\begin{lemma}[Fixed curvature, variable holonomy]\label{lem:tuning}
Let $M$ carry a closed one-form $\eta$ and a loop $\ell$ satisfying
\begin{equation}\label{eq:eta}
d\eta=0,\qquad \int_\ell\eta=2\pi.
\end{equation}
For a global $SO(2)$ connection form $\Ad_0$ on the trivial bundle, define
\begin{equation}\label{eq:Aa}
\Ad_a=\Ad_0+a\eta,\qquad a\in\R.
\end{equation}
Then
\begin{equation}\label{eq:tuning}
d\Ad_a=d\Ad_0,\qquad
\Hol_{\Ad_a}(\ell)=\Hol_{\Ad_0}(\ell)e^{-2\pi aC_1}.
\end{equation}
\end{lemma}

\begin{proof}
The curvature statement follows from $d\eta=0$.  The holonomy statement
follows from Eq.~\eqref{eq:pathtransport}, commutativity of $SO(2)$, and the
normalization in Eq.~\eqref{eq:eta}.
\end{proof}

When $\int_\ell\Ad_0=0$, varying $a$ reaches every target element of $SO(2)$
without changing curvature.

\section{An explicit Lorentzian witness}

Let
\begin{equation}\label{eq:spacetime}
M=\R_t\times S^1_\vartheta\times\R_x\times\R_y,
\end{equation}
and let $\eta$ be the standard closed one-form on $S^1$, normalized by
$\int_{S^1}\eta=2\pi$.  Fix $c,\rho>0$ and define
\begin{equation}\label{eq:metrics}
g=-c^2dt^2+\rho^2\eta^2+dx^2+dy^2,
\qquad
h=c^2dt^2+\rho^2\eta^2+dx^2+dy^2.
\end{equation}
Take the trivial bundle $E=M\times V$.  Let $f\in C^2(\R)$ satisfy
$f'\not\equiv0$, put
\begin{equation}\label{eq:nullcoordinate}
u=t-\frac{x}{c},
\end{equation}
and define
\begin{equation}\label{eq:explicitconnection}
\Ad_a=a\eta+f(u)\,dy.
\end{equation}
Its curvature is independent of $a$:
\begin{equation}\label{eq:explicitF}
F=d\Ad_a=f'(u)\,du\wedge dy.
\end{equation}
Clearly $dF=0$.  The nonzero coordinate components are
\begin{equation}\label{eq:Fcomponents}
F_{ty}=f'(u),\qquad F_{xy}=-\frac1c f'(u).
\end{equation}
Since the metric coefficients are constant,
\begin{equation}\label{eq:maxwellcheck}
\nabla_\mu F^{\mu y}
=\partial_tF^{ty}+\partial_xF^{xy}
=-\frac1{c^2}f''(u)+\frac1{c^2}f''(u)=0,
\end{equation}
and the other divergences vanish.  Hence, in the standard differential-form
description of source-free electromagnetism~\cite{Wald},
\begin{equation}\label{eq:maxwell}
dF=0,\qquad d\star_gF=0,\qquad F\ne0.
\end{equation}
The dependence on $u=t-x/c$ is a propagating plane-wave profile.

Restrict the global bundle, connection, and any global section $q$ to a finite
trivializing cover, using identity transition maps.  Every defect in
Eqs.~\eqref{eq:qdefect}--\eqref{eq:Udefect} vanishes, so
\begin{equation}\label{eq:explicitgluing}
\cF_\Lambda=0.
\end{equation}

Fix
\begin{equation}\label{eq:endpoints}
p=(0,[0],0,0),\qquad q=(T,[0],0,0),
\qquad T>\frac{2\pi\rho}{c},
\end{equation}
and define, for $0\le t\le T$,
\begin{equation}\label{eq:paths}
\gamma_0(t)=(t,[0],0,0),\qquad
\gamma_1(t)=\left(t,\left[\frac{2\pi t}{T}\right],0,0\right).
\end{equation}
Both paths are timelike and join $p$ to $q$.  Their proper times are
\begin{equation}\label{eq:times}
\tau[\gamma_0]=T,
\qquad
\tau[\gamma_1]=T\sqrt{1-\left(\frac{2\pi\rho}{cT}\right)^2},
\end{equation}
hence
\begin{equation}\label{eq:timedifference}
\Delta_{\mathrm{time}}(\gamma_0,\gamma_1)\ne0.
\end{equation}

Let $\ell=\gamma_1*\overline{\gamma_0}$.  Along both paths $dy=0$, while
$\ell$ winds once around $S^1$.  Therefore
\begin{equation}\label{eq:loopholonomy}
\int_\ell f(u)\,dy=0,\qquad
\int_\ell\eta=2\pi,\qquad
\Hol_{\Ad_a}(\ell)=e^{-2\pi aC_1}.
\end{equation}
Choose $a\notin\mathbb Z$.  Then
\begin{equation}\label{eq:phasedifference}
\Delta_{\mathrm{phase}}(\gamma_1,\gamma_0)\ne0.
\end{equation}
Varying $a$ prescribes the holonomy arbitrarily in $SO(2)$ while leaving the
nonzero Maxwell curvature in Eq.~\eqref{eq:explicitF} unchanged.

\begin{proof}[Proof of Theorem~\ref{thm:main}]
Equations~\eqref{eq:generators}--\eqref{eq:dinfty} provide the common carrier,
square symmetry, exact projective factorization, maximal equalizer, and
infinite-dihedral dynamics.  Proposition~\ref{prop:fibpowers},
Eqs.~\eqref{eq:binet}--\eqref{eq:hurwitzlimit}, and
Proposition~\ref{prop:glide} give Fibonacci iteration, the sharp approximation
scale, and the hyperbolic realization.  Equations~\eqref{eq:xiJ} and
\eqref{eq:xiroots} place the completed-zeta reflection without extending its
scope.  Proposition~\ref{prop:gluing} supplies exact local compatibility.
The explicit data in Eqs.~\eqref{eq:spacetime}--\eqref{eq:phasedifference}
satisfy the source-free Maxwell equations, exact gluing, nontrivial relative
holonomy, and unequal proper time simultaneously.  Lemma~\ref{lem:tuning}
proves independent curvature and holonomy tuning.
\end{proof}

\section{The strict boundary}

The construction contains several nonzero or exactly vanishing quantities,
but they are not of one type:
\begin{equation}\label{eq:types}
\begin{aligned}
\epsilon_\Phi &\in\R,\\
F &\in\Omega^2(M),\\
\Delta_{\mathrm{phase}} &\in\operatorname{End}(E_p),\\
\Delta_{\mathrm{time}} &\in\R,\\
\cF_\Lambda &\in\R_{\ge0}.
\end{aligned}
\end{equation}
No equation in this paper identifies these objects or makes one a function of
another.  In particular:
\begin{enumerate}[label=(\alph*),leftmargin=2.1em,itemsep=0.3ex,topsep=0.6ex]
\item the completed-zeta reflection supplies no statement about the location
      of zeta zeros;
\item $M_J\notin O(2)$, so the arithmetic reflection does not act as a gauge
      transformation or fiber isometry;
\item the common generator $C_1$ does not constrain the scalar Maxwell profile
      $f'(u)$;
\item exact gluing $\cF_\Lambda=0$ does not force zero curvature, zero holonomy,
      or equal proper time;
\item nonzero holonomy and unequal proper time are jointly readable but are not
      equated.
\end{enumerate}

This boundary is part of the theorem.  The result is coexistence under exact
typing, not derivation of one canonical theory from another.

\section{Conclusion}

The factorization
\begin{equation}\label{eq:closingfactorization}
M_JC_1=
\begin{pmatrix}1&1\\1&0\end{pmatrix}
\end{equation}
is the smallest algebraic hinge in the construction.  Its alternating
projective action produces the Fibonacci matrices exactly, locates a maximal
two-point meeting of distinct involutions, and fixes a hyperbolic axis whose
length scale is governed by $\Phi$.  The same two-dimensional carrier admits a
separate $SO(2)$ gauge office on a Lorentzian manifold, where exact gluing,
nonzero Maxwell curvature, variable holonomy, and unequal proper time coexist
in one explicit witness.

The point is exact and limited: these structures can inhabit one realization
without being identified.  Compatibility is proved by construction.  Every
stronger bridge remains open.

\section*{Acknowledgments}

The author used generative language-model tools for editorial assistance and
algebraic cross-checking.  Every definition, argument, citation, and conclusion
was reviewed by the author, who accepts full responsibility for the manuscript.

\section*{Author declarations}

\noindent\textit{Conflict of interest.}
The author has no conflicts to disclose.

\medskip
\noindent\textit{Author contributions.}
Anthony Vito Coppa: Conceptualization; Formal analysis; Investigation;
Methodology; Writing, original draft; Writing, review and editing.

\medskip
\noindent\textit{Data availability.}
Data sharing is not applicable to this article because no new data were created
or analyzed.  All mathematical data supporting the result are contained in the
article.

\begin{thebibliography}{99}
\small

\bibitem{BaezMuniain}
J. C. Baez and J. P. Muniain,
\textit{Gauge Fields, Knots and Gravity}
(World Scientific, Singapore, 1994).

\bibitem{Beardon}
A. F. Beardon,
\textit{The Geometry of Discrete Groups}, Graduate Texts in Mathematics,
Vol. 91 (Springer, New York, 1983).

\bibitem{HardyWright}
G. H. Hardy and E. M. Wright,
\textit{An Introduction to the Theory of Numbers}, 6th ed., revised by
D. R. Heath-Brown and J. H. Silverman (Oxford University Press, Oxford, 2008).

\bibitem{Khinchin}
A. Ya. Khinchin,
\textit{Continued Fractions} (Dover, Mineola, NY, 1997).

\bibitem{KobayashiNomizu}
S. Kobayashi and K. Nomizu,
\textit{Foundations of Differential Geometry}, Vol. I
(Wiley, New York, 1963).

\bibitem{Lang}
S. Lang,
\textit{$SL_2(\mathbb R)$}, Graduate Texts in Mathematics, Vol. 105
(Springer, New York, 1985).

\bibitem{Nakahara}
M. Nakahara,
\textit{Geometry, Topology and Physics}, 2nd ed.
(Institute of Physics Publishing, Bristol, 2003).

\bibitem{Titchmarsh}
E. C. Titchmarsh,
\textit{The Theory of the Riemann Zeta-Function}, 2nd ed., revised by
D. R. Heath-Brown (Oxford University Press, Oxford, 1986).

\bibitem{Wald}
R. M. Wald,
\textit{General Relativity} (University of Chicago Press, Chicago, 1984).

\end{thebibliography}

\end{document}
